\section{An honest countable weighted realization}
\label{sec:countable}

The intervalwise incidence identities are algebraic. To form Fredholm
determinants we now realize the selected atom family on a Hilbert space and
prove the required operator ideals directly.

\begin{definition}[Weighted source space]
\label{def:weighted-space}
For \(\eta>\half\), let
\begin{equation}
H_\eta=
\left\{x=(x_n)_{n\ge1}:
\norm{x}_\eta^2=\sum_{n\ge1}n^{2\eta}|x_n|^2<\infty\right\}.
\label{eq:weighted-space}
\end{equation}
\end{definition}

For a source atom \(p\), the kernel \eqref{eq:kernel} gives the explicit
action
\begin{equation}
q_px
=\left(\sum_{k\ge1}\mob(k)x_{pk}\right)(e_1+e_p).
\label{eq:rank-one-action}
\end{equation}
Thus the incidence idempotent is not merely an intervalwise formal element:
it is a concrete rank-one operator.

\begin{theorem}[Rank-one realization and exact trace norm]
\label{thm:rank-one}
For every \(\eta>\half\) and source atom \(p\), the operator \(q_p\) belongs
to \(\Sone(H_\eta)\), satisfies
\[
q_p^2=q_p,\qquad \Tr q_p=1,
\]
and has trace norm
\begin{equation}
\norm{q_p}_1
=\sqrt{(1+p^{-2\eta})C_\eta},
\qquad
C_\eta
=\sum_{k\ge1}\frac{\mob(k)^2}{k^{2\eta}}
=\frac{\zeta(2\eta)}{\zeta(4\eta)}.
\label{eq:trace-norm}
\end{equation}
In particular,
\begin{equation}
\norm{q_p}_1\le\sqrt{2C_\eta}
\label{eq:uniform-trace-norm}
\end{equation}
uniformly in \(p\), and \(q_pq_q=0\) for distinct source atoms.
\end{theorem}

\begin{proof}
Write \eqref{eq:rank-one-action} as
\[
q_p=|r_p\rangle\langle v_p|,
\qquad
r_p=e_1+e_p.
\]
The range vector has
\[
\norm{r_p}_\eta^2=1+p^{2\eta}.
\]
The squared dual norm of the functional is
\[
\norm{v_p}_\eta^2
=\sum_{k\ge1}\frac{\mob(k)^2}{(pk)^{2\eta}}
=p^{-2\eta}C_\eta.
\]
The trace norm of a rank-one operator is the product of these norms, giving
\eqref{eq:trace-norm}. The Euler product for the squarefree indicator gives
\[
C_\eta
=\prod_{\ell\in\Pp}(1+\ell^{-2\eta})
=\frac{\zeta(2\eta)}{\zeta(4\eta)}.
\]
The functional in \eqref{eq:rank-one-action} evaluates to one on
\(e_1+e_p\), proving idempotence and trace one. If \(p\ne q\), neither
\(1\) nor \(q\) is a positive multiple of \(p\); the cross evaluation
vanishes and hence \(q_pq_q=0\).
\end{proof}

The condition \(\eta>\half\) is exactly what is needed for
\(C_\eta<\infty\). The uniform estimate
\eqref{eq:uniform-trace-norm} will control the infinite marked transfer.

\subsection{Bounded similarity of the whole family}

Individual trace-class projectors do not automatically imply a bounded
global change of basis. The stronger similarity statement has a narrower
domain.

\begin{theorem}[Bounded global incidence similarity]
\label{thm:global-similarity}
For \(\eta>1\), the incidence zeta and Möbius transforms extend to bounded
mutually inverse operators \(Z_\eta,M_\eta\) on \(H_\eta\), and
\begin{equation}
q_p=Z_\eta\epscoord_pM_\eta.
\label{eq:countable-similarity}
\end{equation}
Their norms satisfy
\begin{equation}
\norm{Z_\eta}\le\zeta(\eta),
\qquad
\norm{M_\eta}
\le\sum_{k\ge1}\frac{|\mob(k)|}{k^\eta}
=\frac{\zeta(\eta)}{\zeta(2\eta)}.
\label{eq:similarity-bounds}
\end{equation}
\end{theorem}

\begin{proof}
Apply the unitary reweighting \(U_\eta x=(n^\eta x_n)_n\). In unweighted
\(\ell^2\) coordinates,
\begin{equation}
(U_\eta Z_\eta U_\eta^{-1}X)_a
=\sum_{k\ge1}k^{-\eta}X_{ak},
\qquad
(U_\eta M_\eta U_\eta^{-1}X)_a
=\sum_{k\ge1}\mob(k)k^{-\eta}X_{ak}.
\label{eq:downsampling-series}
\end{equation}
Every downsampling map \(X\mapsto(X_{ak})_a\) has norm at most one. The two
operator series therefore converge absolutely in norm for \(\eta>1\), with
the bounds in \eqref{eq:similarity-bounds}. On finitely supported vectors the
incidence identities give \(Z_\eta M_\eta=M_\eta Z_\eta=I\); continuity and
density extend them to all of \(H_\eta\). Equation
\eqref{eq:countable-similarity} follows from the finite-support kernel
identity.
\end{proof}

\begin{remark}[The intermediate strip]
\label{rem:intermediate-eta}
For \(\half<\eta\le1\), every selected \(q_p\) remains an honest rank-one
trace-class idempotent and pair annihilation remains exact. A bounded global
zeta similarity is neither needed for the determinant theorem nor asserted
there. Choosing any \(\eta>1\), however, leaves every symbolic coefficient
unchanged and proves that the canonical whole family can be put in coordinate
form by a bounded similarity.
\end{remark}

The countable theorem therefore sharpens the finite no-go. Oblique incidence
geometry is not a cutoff artifact, but neither does it supply a new ordinary
cyclic character. On a perfectly legitimate common space it is a bounded
change of basis from the coordinate atom family.
